\chapter{Lecture extensions: coverage unconfirmed}
These topics occur in your lectures but are not explicitly asked on the 2025
sample. Attempt after the core sets, or when an announcement confirms scope.
These are supplementary questions, not predictions of exam questions.

\task{9.1. Transformer runner types}
Use these supplied representations:
\par\noindent\begin{minipage}{\linewidth}
\begin{lstlisting}
runStateT  :: StateT s m a -> s -> m (a,s)
runExceptT :: ExceptT e m a -> m (Either e a)
\end{lstlisting}
\end{minipage}\par
Expand each into a function type with no remaining transformer wrappers:
(a) \code{StateT Int (Either String) Bool};
(b) \code{ExceptT String (State Int) Bool}, after running both layers.
Which ordering can return a final state together with an error? Justify from
the expanded result shapes.
\answer{[Expanded types and reasoning.]}

\task{9.2. Immutable arrays and dynamic programming}
Using the Dynamic lecture's array-based edit-distance definition, identify
the base cells and the three predecessor cells for a mismatch. Explain why
calling the naive recursive function rather than looking up the array in the
recursive cases would lose sharing of subproblems.
\answer{[Recurrence, base cases, sharing explanation. Source: Dynamic.hs, Recursion over indices and self-referential array.]}

\clearpage
\task{9.3. GADT evaluator}
Use this indexed expression language. Implement the evaluator directly below.
\par\noindent\begin{minipage}{\linewidth}
\begin{lstlisting}
data Expr a where
  Number :: Int -> Expr Int
  Truth  :: Bool -> Expr Bool
  Plus   :: Expr Int -> Expr Int -> Expr Int
  If     :: Expr Bool -> Expr a -> Expr a -> Expr a
\end{lstlisting}
\end{minipage}\par
\par\noindent\begin{minipage}{\linewidth}
\begin{codeanswer}
evalExpr :: Expr a -> a
evalExpr = undefined  -- YOUR CODE
\end{codeanswer}
\end{minipage}\par
Does the following expression type-check? Explain using the constructor types.
\par\noindent\begin{minipage}{\linewidth}
\begin{lstlisting}
Plus (Truth True) (Number 2)
\end{lstlisting}
\end{minipage}\par
Source: GADTs.hs, A Typed Expression Evaluator.
\answer{[Answer and type conflict, if any.]}

\task{9.4. Agda and propositions as types}
Use the tutorial's correspondence between implication and functions, conjunction
and products, disjunction and sums, and falsity and an empty type.
State the type representing $(P\land Q)\Rightarrow Q$, then implement its proof.
Use the tutorial's product constructor/notation or define an equivalent record.
Explain why negation of $P$ corresponds to $P\to\bot$, rather than to $\bot$ alone.
\par\noindent\begin{minipage}{\linewidth}
\begin{codeanswer}
-- YOUR AGDA TYPE AND PROOF
\end{codeanswer}
\end{minipage}\par
\answer{[Negation explanation. Source: Tutorial.agda, Curry-Howard correspondence.]}
